%% =============================================================================
%% section1.tex — Section 1: Introduction
%% =============================================================================

\section{Introduction}
\label{sec:intro}

% ---- Paragraph 1: Data phenomenon and canonical example ----

Bounded discrete proportions---counts out of known totals---arise
whenever a finite capacity is partially utilized: hospital beds
occupied, insurance plans enrolled, dental surfaces decayed, survey
items endorsed, or enrollment slots filled at a childcare center.  When the outcome can take the value zero for
structural rather than sampling reasons and displays variance well
beyond the binomial, standard generalized linear models are inadequate.
Infant and toddler childcare enrollment in the United States offers a
canonical instance of all these complications simultaneously.  Data from
the 2019 National Survey of Early Care and Education
\citep{NSECEProjectTeam2022} record, for each center-based provider,
the number of children under age three enrolled out of total capacity.
Over one-third of centers report zero infant and toddler enrollment---not
because demand is absent, but because stringent staff-to-child ratios and
elevated per-child costs lead many providers to decline serving this age
group entirely~\citep{AdamsRohacek2010}.  The resulting access gaps fall
disproportionately on high-poverty
communities~\citep{DobbinsEtAl2016, MalikEtAl2018}, yet the relationship
between poverty and enrollment turns out to be considerably more complex
than a simple deficit story.

% ---- Paragraph 2: The poverty reversal → methodological challenge ----

A preliminary examination of the data reveals a distinctive pattern that we
term the \emph{poverty reversal}: community poverty reduces the
probability that a center serves infants and toddlers at all, yet among
centers that \emph{do} participate, those in higher-poverty areas devote
a \emph{larger} share of capacity to this age group.
The reversal poses a methodological challenge: it \emph{cannot be detected} by any model
that conflates participation and intensity into a single equation, nor
by one that fails to accommodate bounded support, overdispersion,
structural zeros, and survey weighting simultaneously.  These four
features jointly motivate the hierarchical hurdle beta-binomial
framework developed in this paper.

% ---- Paragraph 3: Literature review — hurdle/ZI and beta-binomial ----

Two-part models for count data with excess zeros have a substantial
history.  \citet{Cragg1971} introduced the hurdle framework for
continuous expenditures and \citet{Mullahy1986} adapted it for counts,
separating a binary participation stage from a truncated count stage.
Bayesian extensions have expanded their scope considerably: hierarchical
zero-inflated or hurdle models with Poisson and negative-binomial kernels
now accommodate spatial, longitudinal, and multilevel structures
\citep{NeelonEtAl2010, NeelonEtAl2013, NeelonEtAl2016, Neelon2019,
Altinisik2023}.  Most relevant to our work, \citet{GhosalEtAl2020}
propose a hierarchical hurdle model for spatiotemporal count data that
captures cross-margin dependence through a single shared scalar.  However,
all of these models assume Poisson or negative-binomial kernels, which
treat the count as unbounded.  When the response has bounded support, the
beta-binomial distribution is the natural alternative, introducing
overdispersion while respecting the $\{0,1,\ldots,n\}$
support~\citep{WagnerEtAl2015}.  \citet{BandyopadhyayEtAl2011} embed the
beta-binomial in a spatial framework for clustered bounded counts, and
\citet{WenEtAl2025} propose a zero-inflated beta-binomial with spatially
varying coefficients and cross-component covariance.  Yet they adopt a
zero-inflation formulation in which zeros arise from a latent mixture,
whereas in the childcare setting a center reporting zero infant enrollment
has made a deliberate operational decision~\citep[cf.][]{Lambert1992,
Hall2000}, making the hurdle specification the natural choice.

% ---- Paragraph 4: Literature — survey weighting and the gap ----

A separate methodological strand addresses design-based inference in
Bayesian models for survey data.  \citet{SavitskyToth2016} develop a
pseudo-posterior framework in which survey weights enter the likelihood
exponent, and \citet{WilliamsSavitsky2021} provide calibrated sandwich
variance corrections that restore nominal credible-interval
coverage~\citep[building on][]{Binder1983, Pfeffermann1993}.  These
contributions target standard single-equation regressions; their
extension to two-part hierarchical specifications has not been attempted.
In sum, the literature provides the individual building
blocks---hurdle, beta-binomial, hierarchical structure, survey
correction---but not their synthesis.

% ---- Paragraphs 5--6: Two contributions + analytical framework ----

This paper develops that synthesis.  The model extends the framework of
\citet{GhosalEtAl2020} in two directions, each addressing a limitation
of existing approaches.

First, we introduce a Bayesian hierarchical hurdle beta-binomial
model with state-varying coefficients and cross-margin covariance.  The
hurdle component separates participation from intensity; the
beta-binomial kernel respects the bounded support and accommodates the
observed overdispersion, which exceeds the binomial variance by a factor
of approximately twelve; and the hierarchical prior on state-level
coefficients identifies a cross-margin covariance structure whose
effective sample size is $S = 51$ states rather than $N = 6{,}785$
providers, with the $\LKJ$ prior providing finite-sample regularization
for the high-dimensional covariance ($2q = 10$ parameters).  The model nests
the scalar cross-margin parameter of \citet{GhosalEtAl2020} as a special
case and extends the beta-binomial models of
\citet{BandyopadhyayEtAl2011} and \citet{WenEtAl2025} by incorporating a
hurdle for structural zeros.  Applied to the NSECE data, this
specification reveals the reversal---and its geographic
heterogeneity across 51 states---with full uncertainty quantification.

Second, we extend the pseudo-posterior framework of
\citet{SavitskyToth2016} and \citet{WilliamsSavitsky2021} to the
two-part hierarchical setting, implementing a Cholesky-based sandwich
variance correction~\citep{Binder1983, Pfeffermann1993} and introducing
the design effect ratio as a parameter-specific diagnostic.  A companion
simulation study (\cref{sec:simulation}) provides candid diagnostics:
sandwich-corrected intervals substantially improve coverage for fixed
effects under realistic design effects (82--88.5\% at the 90\% nominal
level), while variance components require
unweighted estimation---a finding we report transparently rather than
treating the sandwich as a universal remedy.

These two innovations jointly enable a marginal effect
decomposition~(\cref{prop:lae-lie}) that separates the reversal into
extensive-margin (access) and intensive-margin (intensity) components.
A single-equation specification conflates these offsetting forces into
a muted net effect; the two-part decomposition quantifies each channel
and their relative magnitudes, with the extensive margin accounting for
approximately two-thirds of the total effect.  The companion R package
\textsf{hurdlebb}~\citep{hurdlebb2026} implements the full methodology.

% ---- Paragraph 7: Paper outline ----

The remainder of the paper is organized as follows.
\Cref{sec:data} introduces the NSECE data and documents the empirical
patterns motivating the model.
\Cref{sec:model} develops the hierarchical hurdle beta-binomial
specification, its analytic properties, and the sandwich correction.
\Cref{sec:simulation} reports the simulation study.
\Cref{sec:application} applies the model to the 2019 NSECE data.
\Cref{sec:discussion} discusses implications and limitations.
Supplementary Materials~A--E provide proofs, identification conditions,
extended tables, and computational details.
